\section{Conclusions}
\label{sec:conclusion}

In this work, we studied simplified token mixing modules for Transformer-like encoder architectures, making several contributions. First, we showed that simple, linear mixing transformations, along with the nonlinearities in feed-forward layers, can competently model diverse semantic relationships in text. Second, we introduced FNet, a Transformer-like model wherein the self-attention sublayer is replaced by an unparameterized Fourier Transform. FNets achieve $92$ and $97\%$ of their respective BERT-Base and BERT-Large counterparts' accuracy on the GLUE benchmark, but train $70-80\%$ faster on GPUs/TPUs. Third, because of its favorable scaling properties, FNet is very competitive with the ``efficient Transformers'' evaluated on the Long-Range Arena benchmark, matching the accuracy of the most accurate models while being much faster and lighter on memory. 

Our work highlights the potential of linear units as a drop-in replacement for the attention mechanism in text classification tasks. We found the Fourier Transform to be a particularly efficient and effective mixing mechanism, due to the speed of the FFT. However, we only performed a cursory survey of other linear transformations (see also Appendix \ref{app:things_we_tried}), and additional fast alternatives are worth exploring.

Given the speed and accuracy advantages of smaller FNet models relative to Transformers, we suspect that FNet will be effective as a lightweight, distilled student model deployed in resource-constrained settings such as production services or on edge devices. The need for such lightweight serving models is only forecast to grow given the interest in giant models \citep{raffel2019exploring, brown2020language, lepikhin2020gshard}. A natural avenue to explore in this regard is knowledge distillation of small FNet models from larger Transformer teacher models, following, for example, \citet{sanh2019distilbert, jiao2019tinybert, turc2019well}.

Another aspect of interest and worthy of further study is hybrid FNet-attention models. We found that adding only a few self-attention sublayers to FNet offers a simple way to trade speed for accuracy. Specifically, replacing the final two Fourier sublayers with self-attention provided $97-99\%$ of BERT's accuracy with limited speed penalties.

Throughout this work we have restricted our focus to encoders. FNet decoders can be designed by ``causally'' masking the Vandermonde matrix, but a lower level implementation is required to introduce causal masking to FFTs. How to adapt Fourier mixing for encoder-decoder cross-attention is an open question as evidence suggests that cross-attention may be crucial to performance \citep{you2020hard}. We have focused on tasks which do not require generation so we leave FNet decoders and encoder-decoder setups to future work; although we do remark that the FNet encoder could be used as a drop in replacement in a Transformer as other works have successfully demonstrated; see, for example, \cite{zaheer2020big, guo2021longt5}.
